

\chapterquote{%
  A checkpoint is not a backup.  A backup stores data.
  A checkpoint stores intent --- the accumulated understanding that
  makes the next action meaningful.  Losing a checkpoint is not a
  data-loss event; it is a goal-loss event.%
}{Engineering AI Agents for Production, Volume~1}

\begin{learningbox}
After completing this chapter you will be able to:
\begin{itemize}
  \item Upgrade the BDI fields \code{beliefs} and \code{goals} from
        plain string lists to typed S-P-O triples and hierarchical goal
        nodes, each with its own merge-safe state reducer.
  \item Select the correct LangGraph checkpointer for each deployment
        context and wire \code{AsyncPostgresSaver} with proper connection
        pooling.
  \item Choose between checkpoint-every-cycle and milestone strategies
        based on the cost and recovery semantics of the workflow.
  \item Stamp every checkpoint with a \term{SHA-256 integrity hash} and
        verify it before any state restore, so corrupted or tampered
        snapshots are detected before they can cause goal drift.
  \item Implement a \term{five-step RecoveryNode} that resumes from the
        last valid checkpoint without goal substitution or stale context.
  \item Version the state schema and apply a \term{migration chain} so
        that field additions never break existing checkpoints.
\end{itemize}
\end{learningbox}

% ─────────────────────────────────────────────────────────────────────────────
\section{The BDI State Model}
\label{sec:bdi-model}
% ─────────────────────────────────────────────────────────────────────────────

Chapter~1 declared \code{beliefs} and \code{goals} as lists inside
\code{AgentState}.  This section upgrades their element types to the rich
structures that Volume~1 Chapter~9 specifies.

\subsection{Beliefs as subject--predicate--object triples}
\label{subsec:beliefs}

A free-text string in the beliefs list cannot be deduplicated, queried by
subject, or assigned a confidence score without parsing.  The
\component{Belief} model (\S\ref{subsec:bdi-fields}) replaces the string
with a structured triple.  The merge reducer that LangGraph calls when two
graph branches both update \code{beliefs} must therefore compare by SPO key,
not by string identity:

\begin{lstlisting}[language=Python,
    caption={Merge reducer for beliefs --- dedup by S-P-O key (src/state.py)},
    label={lst:merge-beliefs}]
def _merge_beliefs(
    existing: list[Belief],
    update:   list[Belief],
) -> list[Belief]:
    """
    Append beliefs from update that are not already in existing.
    Deduplication key: (subject, predicate, object_).
    A higher-confidence assertion of the same fact replaces the old one
    at the next write cycle; it does not create a duplicate entry here.
    """
    existing_keys = {
        (b.subject, b.predicate, b.object_)
        for b in existing
    }
    return existing + [
        b for b in update
        if (b.subject, b.predicate, b.object_) not in existing_keys
    ]
\end{lstlisting}
\githubref{src.state.\_merge\_beliefs}

\medskip
The reducer is registered in the \code{AgentState} TypedDict with the
\code{Annotated} syntax so LangGraph calls it automatically on every
state merge:

\begin{lstlisting}[language=Python,
    caption={Reducer annotation in AgentState (src/state.py)},
    label={lst:belief-annotation}]
class AgentState(TypedDict, total=False):
    beliefs: Annotated[list[Belief], _merge_beliefs]
    goals:   Annotated[list[Goal],   _merge_goals]
\end{lstlisting}

\subsection{Goals as a hierarchical stack}
\label{subsec:goals}

The \code{Goal} model carries a \code{parent\_id} field that links
sub-goals to their parent, an \code{allocated\_budget} for token
sub-allocation across a planning hierarchy (Chapter~15), and
\code{excluded\_tools} that the semantic action validator (Chapter~1
\S\ref{sec:action-node}) checks before every dispatch.

The goal merge reducer resolves conflicts by id --- later writes to the
same goal id win, which is the correct behaviour when the review subgraph
marks a sub-goal \code{COMPLETED} or \code{FAILED}:

\begin{lstlisting}[language=Python,
    caption={Merge reducer for goals --- update by id (src/state.py)},
    label={lst:merge-goals}]
def _merge_goals(
    existing: list[Goal],
    update:   list[Goal],
) -> list[Goal]:
    """
    Update existing goals by id; append genuinely new ones.
    Later writes to the same id override the earlier entry,
    so a COMPLETED status propagates correctly from the review subgraph.
    """
    by_id = {g.id: g for g in existing}
    for g in update:
        by_id[g.id] = g          # last write wins per id
    return list(by_id.values())
\end{lstlisting}
\githubref{src.state.\_merge\_goals}

\begin{tipbox}
\textbf{Every list field in \code{AgentState} needs an explicit reducer.}
Without one, LangGraph replaces the entire list rather than merging.
This is safe for \code{messages} (which has the built-in
\code{add\_messages} reducer that deduplicates by message id) but
dangerous for \code{beliefs} and \code{goals}, where replacement would
lose beliefs written by a parallel branch.
\end{tipbox}

\subsection{Working memory tiers}
\label{subsec:wm-tier}

The \code{working\_memory} list uses an update-by-id reducer
(\code{\_merge\_wm}) and its entries carry a \code{WorkingMemoryTier}
that the eviction logic in Chapter~7 reads.  The four tiers, in ascending
eviction priority, are:

\begin{table}[htbp]
\centering
\caption{WorkingMemoryTier eviction order (lower priority evicted first)}
\label{tab:wm-tiers}
\renewcommand{\arraystretch}{1.4}
\begin{tabularx}{\linewidth}{@{} l l P{6.0cm} @{}}
\toprule
\textbf{Priority} & \textbf{Tier} & \textbf{Used for} \\
\midrule
0 (first out) & \code{ARCHIVED}    & Superseded observations; safe to evict \\
1             & \code{COMPRESSED}  & Summarised tool error output \\
2             & \code{BACKGROUND}  & Routine tool results; context \\
3 (last out)  & \code{ACTIVE}      & Disambiguation entries from recovery nodes;
                                     never evicted \\
\bottomrule
\end{tabularx}
\end{table}

% ─────────────────────────────────────────────────────────────────────────────
\section{Choosing a Checkpointer}
\label{sec:checkpointer}
% ─────────────────────────────────────────────────────────────────────────────

LangGraph ships two async checkpointers.  \component{AsyncSqliteSaver} is
appropriate for development and single-process testing --- it requires no
running infrastructure but cannot serve concurrent graph resumptions from
different processes.  \component{AsyncPostgresSaver} is the production
choice: it serialises checkpoints into a PostgreSQL table, survives process
restarts, and supports concurrent access across multiple worker instances.

\begin{lstlisting}[language=Python,
    caption={make\_checkpointer --- factory for dev and production
             (src/checkpointing.py)},
    label={lst:make-checkpointer}]
from contextlib import asynccontextmanager

@asynccontextmanager
async def make_checkpointer(
    mode:        Literal["sqlite", "postgres"] = "sqlite",
    conn_string: str | None = None,
):
    if mode == "sqlite":
        path = conn_string or ":memory:"
        async with AsyncSqliteSaver.from_conn_string(path) as saver:
            yield saver

    elif mode == "postgres":
        import asyncpg
        from langgraph.checkpoint.postgres.aio import AsyncPostgresSaver

        dsn  = conn_string or os.environ["POSTGRES_DSN"]
        pool = await asyncpg.create_pool(dsn, min_size=2, max_size=10)
        saver = AsyncPostgresSaver(pool)
        await saver.setup()      # creates checkpoint table if absent
        try:
            yield saver
        finally:
            await pool.close()
\end{lstlisting}
\githubref{src.checkpointing.make\_checkpointer}

\begin{warningbox}
\textbf{Call \code{saver.setup()} before the first graph invocation.}
\code{AsyncPostgresSaver.setup()} creates the \code{checkpoints} table
when it does not exist.  Skipping this call causes an
\code{UndefinedTableError} on the first checkpoint write.  In a
containerised deployment, run \code{setup()} in the application startup
hook before the first request arrives, not lazily inside a request handler.
\end{warningbox}

The \code{@asynccontextmanager} pattern ensures the connection pool is
closed cleanly even when the agent raises an unhandled exception.  Both
modes are consumed identically by the caller:

\begin{lstlisting}[language=Python,
    caption={Consuming make\_checkpointer in application code},
    label={lst:use-checkpointer}]
async def main() -> None:
    async with make_checkpointer("postgres") as cp:
        graph = build_agent_graph(checkpointer=cp, policy=policy)
        result = await run_agent(graph, state, policy, guard, agent_id)
\end{lstlisting}

% ─────────────────────────────────────────────────────────────────────────────
\section{Checkpoint Strategy}
\label{sec:checkpoint-strategy}
% ─────────────────────────────────────────────────────────────────────────────

LangGraph writes a checkpoint by default after every node execution.
For a four-node graph plus two subgraphs, that is six writes per cycle.
The \component{CheckpointStrategy} class wraps the raw checkpointer
to thin down the write frequency when per-cycle checkpointing is
unacceptable.

\begin{lstlisting}[language=Python,
    caption={CheckpointStrategy --- every-cycle versus milestone
             (src/checkpointing.py)},
    label={lst:checkpoint-strategy}]
class CheckpointStrategy:
    """
    every_cycle: LangGraph default — one checkpoint per node per cycle.
                 Use for: HITL flows, short-lived agents, any run where
                 replay from an arbitrary point is required.
    milestone:   checkpoint only when cycle_count % interval == 0.
                 Use for: long-running background jobs where per-node
                 write latency is unacceptable.
    """
    def __init__(
        self,
        mode:     Literal["every_cycle", "milestone"] = "every_cycle",
        interval: int = 5,
    ) -> None:
        self.mode     = mode
        self.interval = interval

    def should_checkpoint(self, cycle: int) -> bool:
        if self.mode == "every_cycle":
            return True
        return cycle % self.interval == 0
\end{lstlisting}
\githubref{src.checkpointing.CheckpointStrategy}

\medskip
\noindent
\textbf{Which strategy to use.}  The \code{every\_cycle} strategy is
required for HITL flows (Chapter~21): if the graph is interrupted for
human approval, the paused state must be available immediately.  The
\code{milestone} strategy is acceptable for batch jobs (Chapter~22) where
a resume after a crash replays a bounded number of cycles from the last
milestone rather than resuming from the exact interrupted point.

% ─────────────────────────────────────────────────────────────────────────────
\section{Snapshot Integrity with SHA-256}
\label{sec:integrity}
% ─────────────────────────────────────────────────────────────────────────────

Volume~1 §13.4 requires that every checkpoint restore is preceded by an
integrity check.  The Java framework's
\code{AgentRuntime.replay(snapshot, agent)} calls
\code{verifyIntegrity(snapshot)} before applying a single field from the
stored state.  This catches three failure modes: a storage write that was
only partially flushed, a manual database edit during debugging, and ---
in hostile environments --- deliberate tampering.

\begin{lstlisting}[language=Python,
    caption={SHA-256 snapshot integrity --- stamp and verify
             (src/checkpointing.py)},
    label={lst:integrity}]
def compute_snapshot_hash(state: dict) -> str:
    """
    Deterministic SHA-256 over the canonical state dict.
    Excluded fields: _integrity_hash itself, pending_approval
    (session-specific), cycle_traces and tool_records (append-only
    audit lists that change between stamp and verify).
    """
    exclude  = {"_integrity_hash", "pending_approval",
                "cycle_traces",    "tool_records"}
    snapshot = {k: v for k, v in state.items()
                if k not in exclude and v is not None}

    def _serialise(obj):
        if hasattr(obj, "model_dump"):  return obj.model_dump()
        if hasattr(obj, "value"):       return obj.value
        if isinstance(obj, list):       return [_serialise(x) for x in obj]
        if isinstance(obj, dict):
            return {k: _serialise(v) for k, v in sorted(obj.items())}
        return obj

    canonical = json.dumps(_serialise(snapshot), sort_keys=True, default=str)
    return hashlib.sha256(canonical.encode()).hexdigest()

def stamp_integrity(state: dict) -> dict:
    return {**state, "_integrity_hash": compute_snapshot_hash(state)}

def verify_integrity(state: dict) -> bool:
    stored = state.get("_integrity_hash")
    if not stored:
        return True       # legacy checkpoint without hash — permit with log
    return stored == compute_snapshot_hash(state)
\end{lstlisting}
\githubref{src.checkpointing.verify\_integrity}

\medskip
The \code{\_serialise} helper converts Pydantic models to dicts and
\code{Enum} members to their string values before hashing, so the hash is
stable across Python session restarts and does not depend on object
identity.  The four excluded fields are those that legitimately differ
between the moment a checkpoint is stamped and the moment it is verified:
\code{cycle\_traces} and \code{tool\_records} are append-only; a resumed
run will add entries before the first verify call.

\begin{notebox}
\textbf{Stamp at write, verify at read.}  \code{stamp\_integrity()} is
called by the checkpointer wrapper immediately before writing to storage.
\code{verify\_integrity()} is called by \code{RecoveryNode} as the first
step before restoring any field.  The window between stamp and verify is
exactly the storage round-trip.
\end{notebox}

% ─────────────────────────────────────────────────────────────────────────────
\section{The RecoveryNode: Five Steps to Safe Resume}
\label{sec:recovery-node}
% ─────────────────────────────────────────────────────────────────────────────

Recovery from a checkpoint is not a simple state copy.  A naive restore
can re-introduce a goal that was already completed in the lost cycles, or
carry forward a stale working-memory entry that was invalidated by a tool
result that never made it into storage.  Volume~1 §13.4 specifies a
five-step sequence; the \component{RecoveryNode} implements it in order.

\begin{lstlisting}[language=Python,
    caption={\texttt{RecoveryNode} --- five-step safe resume
             (src/checkpointing.py)},
    label={lst:recovery-node}]
@observe(name="recovery_node")
async def recovery_node(
    state:  AgentState,
    config: RunnableConfig,
    graph:  CompiledStateGraph,
) -> dict:
    run_config = config or RunnableConfig(
        configurable={"thread_id": state.get("run_id", "recovery")}
    )

    # Step 1 — Load the last checkpoint from storage
    saved    = await graph.aget_state(run_config)
    snapshot = dict(saved.values)

    # Step 2 — Integrity verification (SHA-256)
    if not verify_integrity(snapshot):
        return {
            "run_state":          RunState.TERMINATED,
            "termination_reason": TerminationReason.ESCALATED,
            "dead_letter": {"error": "Checkpoint integrity check failed"},
        }

    # Step 3 — Schema migration (if snapshot is from an older version)
    from_version = snapshot.get("schema_version", 1)
    if from_version < CURRENT_SCHEMA_VERSION:
        snapshot = migrate_state(snapshot, from_version)

    # Step 4 — Validate structural invariants
    errors = _validate_restored_state(snapshot)
    if errors:
        return {
            "run_state":          RunState.TERMINATED,
            "dead_letter": {"validation_errors": errors},
        }

    # Step 5 — Re-inject original goal; clear per-cycle ephemeral fields
    snapshot["goal"]           = state.get("goal") or snapshot.get("goal","")
    snapshot["decision"]       = None
    snapshot["action_result"]  = None
    snapshot["reflect_output"] = None
    snapshot["run_state"]      = RunState.INITIALIZED

    return snapshot
\end{lstlisting}
\githubref{src.checkpointing.recovery\_node}

\medskip
Step~5 re-injects the \emph{original} goal from the incoming state rather
than from the checkpoint.  This prevents a category of attack called
\term{goal substitution}: if a hostile payload reached the agent in the
lost cycles and modified the goal field before the crash, restoring from
the checkpoint would resume with the corrupted goal.  The incoming state's
goal is the one the operator submitted at the start of the run; it is never
derived from stored data.

The ephemeral fields (\code{decision}, \code{action\_result},
\code{reflect\_output}) are cleared so the agent re-enters at
\code{INITIALIZED} and runs a fresh perception cycle.  Carrying a stale
\code{decision} from a crashed cycle into the next invocation would
dispatch a tool call whose context no longer applies.

% ─────────────────────────────────────────────────────────────────────────────
\section{State Schema Migration}
\label{sec:migration}
% ─────────────────────────────────────────────────────────────────────────────

As the application evolves, \code{AgentState} gains new fields.  A
checkpoint written by version~1 of the code does not have those fields.
The migration chain in \code{src/checkpointing.py} applies transformations
from any stored version up to \code{CURRENT\_SCHEMA\_VERSION}:

\begin{lstlisting}[language=Python,
    caption={Migration chain --- v1 to v2 (src/checkpointing.py)},
    label={lst:migration}]
CURRENT_SCHEMA_VERSION: int = 2

MIGRATION_CHAIN: dict[int, dict] = {
    1: {
        # v1 had beliefs: list[str] and goals: list[str]
        # v2 requires beliefs: list[Belief] and goals: list[Goal]
        "beliefs":        lambda old: _migrate_beliefs_v1(old),
        "goals":          lambda old: _migrate_goals_v1(old),
        "schema_version": lambda _:   2,
    }
}

def migrate_state(state: dict, from_version: int) -> dict:
    current = dict(state)
    version = from_version
    while version < CURRENT_SCHEMA_VERSION:
        migrations = MIGRATION_CHAIN.get(version)
        if not migrations:
            break
        for field, migrate_fn in migrations.items():
            if field in current:
                current[field] = migrate_fn(current[field])
            elif field == "schema_version":
                current["schema_version"] = migrate_fn(None)
        version += 1
    return current
\end{lstlisting}
\githubref{src.checkpointing.migrate\_state}

\medskip
The migration functions for beliefs and goals convert each string element
to the corresponding Pydantic dict representation, preserving the original
text in the \code{object\_} field of a synthetic Belief and in the
\code{description} field of a Goal.  The converted entries carry
\code{provenance = "migration:v1\_to\_v2"} so operators can identify
migrated data in audit queries.

Adding a new field in a future schema version requires only three steps:
declare the field in \code{AgentState} with a default value, increment
\code{CURRENT\_SCHEMA\_VERSION}, and add a new entry to
\code{MIGRATION\_CHAIN} that sets the field's default for older snapshots.
The migration is applied transparently inside \code{recovery\_node}
before any validator sees the restored state.

\begin{tipbox}
\textbf{Migrations must be idempotent.}  \code{migrate\_state()} may be
called more than once on the same checkpoint if a recovery flow is
interrupted and retried.  Each migration function must detect whether the
field is already in the target format and skip the conversion if so.  The
version check \code{while version < CURRENT\_SCHEMA\_VERSION} ensures the
loop terminates, but individual functions must not assume they always
receive the old format.
\end{tipbox}

% ─────────────────────────────────────────────────────────────────────────────
\section{Validating the Checkpointing Layer}
\label{sec:ch3-tests}
% ─────────────────────────────────────────────────────────────────────────────

\begin{lstlisting}[language=Python,
    caption={Representative Chapter 3 unit tests (tests/unit/test\_ch03.py)},
    label={lst:ch03-tests}]
def test_sha256_detects_tampering():
    state   = {"goal": "test", "cycle_count": 1, "schema_version": 2}
    stamped = stamp_integrity(state)
    assert verify_integrity(stamped)
    tampered = {**stamped, "cycle_count": 0}
    assert not verify_integrity(tampered)

def test_cycle_traces_excluded_from_hash():
    """Audit lists grow after stamping; they must not invalidate the hash."""
    s1 = stamp_integrity({"goal": "g", "schema_version": 2})
    s2 = stamp_integrity({**s1, "cycle_traces": ["a", "b", "c"]})
    assert s1["_integrity_hash"] == s2["_integrity_hash"]

def test_migration_v1_beliefs_to_v2():
    v1 = {"goal": "test", "schema_version": 1,
          "beliefs": ["user asked for X", "tool returned Y"]}
    v2 = migrate_state(v1, from_version=1)
    assert v2["schema_version"] == 2
    assert all(isinstance(b, dict) and "subject" in b for b in v2["beliefs"])
    assert all("migration:v1_to_v2" in b["provenance"] for b in v2["beliefs"])

def test_migration_idempotent():
    v1 = {"goal": "test", "schema_version": 1, "beliefs": ["fact A"]}
    v2 = migrate_state(v1, from_version=1)
    v2_again = migrate_state(v2, from_version=2)
    assert v2_again["schema_version"] == 2

def test_merge_beliefs_deduplicates_by_spo():
    b = Belief(subject="s", predicate="p", object_="o", confidence=0.9)
    b_dup = Belief(subject="s", predicate="p", object_="o", confidence=0.5)
    merged = _merge_beliefs([b], [b_dup])
    assert len(merged) == 1       # same SPO — not duplicated

def test_merge_goals_updates_by_id():
    g_active    = Goal(id="g1", description="Find papers", status=GoalStatus.ACTIVE)
    g_completed = Goal(id="g1", description="Find papers", status=GoalStatus.COMPLETED)
    merged = _merge_goals([g_active], [g_completed])
    assert len(merged) == 1
    assert merged[0].status == GoalStatus.COMPLETED
\end{lstlisting}
\githubref{tests.unit.test\_ch03}

% ─────────────────────────────────────────────────────────────────────────────
\begin{takeawaybox}[Chapter~3 Key Takeaways]
\begin{itemize}
  \item \textbf{BDI fields require typed elements and merge reducers.}
        \code{beliefs} is a list of S-P-O \code{Belief} triples
        deduplicated by \code{(subject, predicate, object\_)}.
        \code{goals} is a list of \code{Goal} nodes updated by id.
        Parallel branches that both write to these fields will merge
        correctly without losing entries from either branch.

  \item \textbf{Dev and production checkpointers differ in resilience,
        not in API.}  Both are consumed through the same
        \code{make\_checkpointer()} context manager.  The switch from
        SQLite to PostgreSQL requires only an environment variable change,
        not a code change.

  \item \textbf{Integrity verification is non-negotiable.}
        Every resume path calls \code{verify\_integrity()} before
        reading a single field from the checkpoint.  A mismatch means
        the run terminates with \code{TerminationReason.ESCALATED}, not
        silently resumes with corrupted state.

  \item \textbf{Five ordered steps separate a safe resume from a
        dangerous one.}
        Load → verify → migrate → validate → re-inject-goal-and-clear.
        Skipping the re-injection step exposes the agent to goal
        substitution; skipping the migration step causes type errors on
        newly added fields.

  \item \textbf{Schema migration is additive and idempotent.}
        New fields always have default values.  Migration functions detect
        already-migrated data and skip.  \code{CURRENT\_SCHEMA\_VERSION}
        is the single number that controls which migrations apply.
\end{itemize}
\end{takeawaybox}
